\section*{第 \ref{ch:b3:electrons-in-solids} 章——固体中的电子}

\begin{solution}{exo:b3:electrons-in-solids:1}
(a) $\tau = m_{\text{e}}/ne^2\rho = \qty{2.5e-14}{s}$。(b) $\ell =
v_{\text{F}}\tau \approx \qty{39}{nm}$。(c) 约 110 个\omterm{def:b3:crystalline-solids:lattice}{晶格
常数}。(d) 一个电子穿过一百个离子而
未被散射，说明它不可能是被离子本身弹开的——晶格
必须对它透明，这正是布洛赫定理
（\cref{thm:b3:electrons-in-solids:bloch}）后来证明的结论。
\end{solution}

\begin{solution}{exo:b3:electrons-in-solids:2}
(a) $v_{\text{d}} = I/neA = \qty{7.3e-4}{m/s}$——不到
每秒一毫米。(b) 约 23 分钟。(c) \emph{场}
沿导线以接近光速自行建立，并使整个电子海
同时漂移：行军的人很慢，
而行军的命令并不慢。(d) $v_{\text{d}}/v_{\text{F}} \sim
10^{-9}$：导电是叠加于剧烈量子运动之上的一种不易察觉的偏置，
而非平稳的经典流动。
\end{solution}

\begin{solution}{exo:b3:electrons-in-solids:3}
(a) 半满能带：\omterm{def:b3:electrons-in-solids:classes}{金属}，$\dd\rho/\dd T > 0$。(b) 两
个电子恰好充满该能带——镁之所以能导电，
是因为其满带与下一个空带\emph{交叠}：\omterm{def:b3:electrons-in-solids:classes}{金属}，
$\dd\rho/\dd T > 0$。(c) 满带，\qty{5.5}{eV} 带隙：\omterm{def:b3:electrons-in-solids:classes}{绝缘体}
（形式上 $\dd\rho/\dd T < 0$，但几乎没有载流子
可供计数）。(d) \omterm{def:b3:electrons-in-solids:classes}{半导体}：$\dd\rho/\dd T < 0$，
这一指数式的符号翻转，正暴露了它的类别。
\end{solution}

\begin{solution}{exo:b3:electrons-in-solids:4}
(a) $\lambda = hc/E_{\text{g}} = \qty{1.1}{\micro m}$：硅
在该波长以外的红外区域是透明的——这也是红外
照相机能用硅做透镜的原因。(b) \qty{225}{nm}，位于
紫外区：没有可见光光子能被吸收，因此金刚石
像水一样清澈。(c) \qty{365}{nm}：能容纳蓝光的带隙
（\qty{2.8}{eV}）只存在于氮化物中，几乎
在别处都难以实用——蓝色发光二极管为材料生长等待了数十年。
(d) 带隙决定其光子能量：过驱动增加的是热量，而非
带隙宽度，二极管会发热而变暗——发光二极管的颜色是
材料常数，而非亮度设定。
\end{solution}

\begin{solution}{exo:b3:electrons-in-solids:5}
(a) $\kappa/\sigma T = 400/(5.9\times10^7\times300) =
\qty{2.3e-8}{W.\ohm/K^2}$。(b) 与 $L$ 相差在百分之十以内。(c)
费米面上的同一批电子既携带电荷，也携带
热量，因此它们的散射时间在比值中相消。(d)
钢锅能把热量传给食物，是因为其中的电子在运动；
木柄在两种意义上都是\omterm{def:b3:electrons-in-solids:classes}{绝缘体}，把电子——
连同热量——一并挡在手之外。
\end{solution}

\begin{solution}{exo:b3:electrons-in-solids:6}
(a) 多出百分之五十——与实验明显不符。
(b) $T/T_{\text{F}} \approx 0.004$：电子项
被抑制到百分之一以下。(c) 晶格的 $T^3$ 项
衰减得比电子的 $T$ 项更快：在几开尔文以下
“可忽略”的电子便是剩下的全部。(d) $C/T =
\gamma + \beta T^2$：截距 $\gamma$ 度量电子的贡献，
斜率 $\beta$ 则反映晶格——一张图，两类住户俱全。
\end{solution}

\begin{solution}{exo:b3:electrons-in-solids:7}
(a) $k_{\text{F}} = (3\pi^2n)^{1/3} = \qty{1.36e10}{m^{-1}}$：
$\lambda_{\text{F}} = \qty{4.6}{\angstrom}$。(b) $2a =
\qty{7.2}{\angstrom}$：量级相同——铜的费米面
延伸向布里渊区边界，而打开能隙的物理过程
正\emph{在}决定性的能量处发生。(c) 一个驻波
把电荷堆积在正离子上（静电能较低），
另一个位于离子之间（能量较高）：波长相同，两种
能量——即能隙。(d) 110 个\omterm{def:b3:crystalline-solids:lattice}{晶格常数}的自由程：
布洛赫波不会被完美晶格散射，因此只剩
振动与杂质对其构成阻碍。
\end{solution}

\begin{solution}{exo:b3:electrons-in-solids:8}
(a) $E_{\text{g}}/2k_{\text{B}} = \qty{6500}{K}$：
$\exp[6500(1/300 - 1/400)] \approx 230$。(b) 热敏电阻的
电阻呈指数下降——曲线陡峭、灵敏；
铂的电阻则平缓、线性地上升（声子更多）。一种是
以载流子数为依据的温度计，另一种以散射为依据。(c)
载流子以电子–\omterm{def:b3:electrons-in-solids:doping}{空穴}\emph{对}的形式产生，而
平衡关系 $np = n_{\text{i}}^2$ 把能隙的代价在
两者之间分摊——因此为 $E_{\text{g}}/2$。(d)
$n_{\text{i}}$ 达到 $\qty{5e21}{m^{-3}}$，温度接近 \qty{760}{K}：
\omterm{def:b3:electrons-in-solids:doping}{掺杂}效应被淹没，电路便丧失了原有的设计。实际
器件会更早失效——其反向漏电流每
升高十开尔文便翻一番，最先失常。
\end{solution}

\begin{solution}{exo:b3:electrons-in-solids:9}
(a) $\qty{5e22}{m^{-3}}$——为 $n_{\text{i}}$ 的五百万倍。
(b) 同样的倍数 $\sim5\times10^6$：每百万个原子中
仅一个杂质原子便完全决定了电导率。(c) 因为意外混入的 ppm
级杂质同样会起作用，且不请自来：只有纯度达 ppb 级的\omterm{def:b3:crystalline-solids:lattice}{晶体}，
其电导率才真正属于设计者。(d) $d \sim n^{-1/3}
\approx \qty{27}{nm}$，是 \qty{2.4}{nm} 轨道的十倍：\omterm{def:b3:electrons-in-solids:doping}{施主}
是孤立的氢原子；把\omterm{def:b3:electrons-in-solids:doping}{掺杂}量提高百倍，轨道便
相接触——形成杂质带，并最终成为\omterm{def:b3:electrons-in-solids:classes}{金属}。
\end{solution}

\begin{solution}{exo:b3:electrons-in-solids:10}
(a) $E = 13.6\times0.26/11.7^2 \approx \qty{26}{meV}$（
磷的实测值 \qty{45}{meV} 使该标度
可信）。(b) 与 $k_{\text{B}}T = \qty{26}{meV}$ 相当：
几乎所有\omterm{def:b3:electrons-in-solids:doping}{施主}在室温下都已电离——
这正是 \cref{def:b3:electrons-in-solids:doping} 的前提。(c) $a =
\qty{0.53}{\angstrom}\times\varepsilon_{\text{r}}/(m^*/
m_{\text{e}}) \approx \qty{2.4}{nm}$。(d) 跨越
数十个\omterm{def:b3:crystalline-solids:lattice}{晶胞}的轨道把\omterm{def:b3:crystalline-solids:lattice}{晶体}视为平滑介质——这
正是使体相的
$\varepsilon_{\text{r}}$ 与能带平均的 $m^*$ 得以成立的条件：
这一近似具有自洽性。
\end{solution}

\begin{solution}{exo:b3:electrons-in-solids:11}
(a) $V_{\text{bi}} = 0.0259\ln(10^{12}) \approx \qty{0.71}{V}$。
(b) $eV/k_{\text{B}}T = 19.3$：比值 $\eu^{19.3} \approx
2\times10^{8}$——二极管是一条精确到八位
数字的单向通道。(c) $I_0 \propto n_{\text{i}}^2 \propto
\eu^{-E_{\text{g}}/k_{\text{B}}T}$：与
\cref{exo:b3:electrons-in-solids:8} 中的指数相同，故有
经验性的加倍规律。(d) 四个二极管构成一个菱形：每个
半周期内，无论哪一对处于正向偏置，都使电流
经负载以\emph{相同}方向流过——交流输入，脉动
直流输出。
\end{solution}

\begin{solution}{exo:b3:electrons-in-solids:12}
(a) 减小带隙，更多光子能够越过带隙，但每个光子
只能提供较小的带隙能量；增大带隙，每个光子
提供的能量更多，但符合条件的光子更少：收集量 $\times$ 电压在两者之间取得峰值。
(b) 略低于 \qty{1.3}{eV} 的最优值：硅的理想
上限约为 $\sim30\,\%$——已足够接近，因而其低廉的价格占了上风。
(c) 宽带隙的顶层电池在高电压下吸收蓝光，并
将红光传给下方的窄带隙电池：每个光子都在
接近其自身能量处被收集，热化损失减小，叠层的
上限升至四十多个百分点。(d) 温度升高使 $I_0$
按指数增大，从而压低开路电压
$\sim(k_{\text{B}}T/e)\ln(I_{\text{L}}/I_0)$；同时带隙本身
略有减小，以电压换取电流，而这部分电流
无法完全挽回。
\end{solution}

\begin{solution}{pb:b3:electrons-in-solids:1}
\textbf{1.} $\lambda_{\text{峰}} \approx \qty{500}{nm}$：
约 \qty{2.5}{eV}。
\textbf{2.} $hc/E_{\text{g}} \approx \qty{1.1}{\micro m}$。
\textbf{3.} 光穿过电池而不被吸收（没有可供吸收进入的终态），
最终以热的形式留在背板中——加热的是屋顶，而非
导线。
\textbf{4.} $\qty{1.12}{eV}$ 以电子-空穴对的形式保留；多余的
\qty{1.4}{eV} 在皮秒量级内经声子耗散——比任何电路
所能截获的都要快。
\textbf{5.} 次带隙透明（约占功率的 $\sim20\,\%$）与
超带隙热化（约占 $\sim30\,\%$）：光谱
在结区见到第一个电子之前就已损失一半。
\textbf{6.} 吸收是单光子量子跃迁，需要一个
真实的终态；在半个带隙处并没有可供停留的态，而
在太阳光的光子密度下，双光子事件可以忽略不计。
\textbf{7.} 在\omterm{prop:b3:electrons-in-solids:junction}{耗尽区}内或其附近产生的电子-空穴对感受到
内建电场：电子被扫至 n 区，\omterm{def:b3:electrons-in-solids:doping}{空穴}被扫至
p 区，电场分开它们的速度快于它们
相互找到并复合的速度——电荷不断积累，直至
外电路将其释放为电流。
\textbf{8.} $V_{\text{bi}} = 0.0259\ln(10^{12}) \approx
\qty{0.71}{V}$。
\textbf{9.} $60\times\qty{0.6}{V} = \qty{36}{V}$。
\textbf{10.} $I = 400/36 \approx \qty{11}{A}$。
\textbf{11.} 整个面板上，超带隙光的功率为
$\sim\qty{1400}{W}$，每个光子能量约 $\sim\qty{2}{eV}$：
每秒 $4\times10^{21}$ 个光子——但六十个电池是\emph{串联}的，
因此 \qty{11}{A} 的电流（电子流 $I/e \approx
7\times10^{19}\,\unit{s^{-1}}$）流经每一个电池，其
各自分得的光子数为 $4\times10^{21}/60 \approx
7\times10^{19}$ 每秒：几乎每个被吸收的光子都产生
一个被收集的电子。量子效率极高；
损失在于能量，而不在于数量。
\textbf{12.} 在光谱的 50\,\% 之后，电池要付出电压不足
（$0.6/1.12 \approx 54\,\%$）以及少量反射与复合的代价：
$0.5\times0.54 \approx 27\,\%$，
修正为额定的 20\,\%。
\textbf{13.} 漏电流 $I_0$：开路电压
$(k_{\text{B}}T/e)\ln(I_{\text{L}}/I_0)$ 停止增长于
光电流与二极管漏电流相平衡之处——约为 \qty{0.6}{V}，远
低于带隙所对应的值。
\textbf{14.} $\sin30^\circ = 0.5$：仅由几何因素就给出正午额定
功率的一半——尚未计入云层的影响。
\textbf{15.} $\Delta T = \qty{40}{K}$：下降 $-16\,\%$，剩下
$\approx\qty{335}{W}$——阳光最充足的日子并不是
每瓦发电量最高的日子。
\textbf{16.} 温度升高使 $I_0$ 按指数增大，而
开路电压——对数关系的“报应”——每开尔文下降
不到百分之一；带隙的轻微缩小
使这一过程进一步加剧。
\textbf{17.} 串联接线迫使所有电池通过同一电流：被遮挡的
电池本身不产生电流，却被其余
五十九个电池驱动进入反向偏置，并以热斑形式耗散它们的功率
——整串电路受限于最弱的那一个电池。
\textbf{18.} 旁路二极管为电流提供一条正向通路，
绕过被遮挡的子串，牺牲的是其电压，而非
整个面板的输出。
\textbf{19.} $\qty{400}{W}\times0.15\times\qty{8760}{h} \approx
\qty{530}{kWh}$ 每年。
\textbf{20.} $\approx\qty{130}{euros}$ 每年：回收期
约 $300/130 \approx 2.5$ 年，此后二十年均为盈利。
\textbf{21.} $500/530 \approx$ 一年：面板回收其
制造能耗的速度，与其价格回收的速度大致相当。
\textbf{22.} 黑色并非一种选择，而是由\omterm{thm:b3:electrons-in-solids:bloch}{能带结构}决定的：
吸收需要一个比满态高出一个光子能量的空态，
而在带隙以下，硅根本没有可提供的态——
涂料无法增加态。
\textbf{23.} 在顶部叠放一个宽带隙电池，以在高电压下吸收蓝光，
并将红光传向下方：该叠层结构突破了
单结的上限。难点在于：串联叠层必须使
各层之间的电流（以及成本）相互匹配。
\textbf{24.} 凡是缩短 $\tau$ 并损害结的因素：
紫外线产生的缺陷与向内扩散的杂质会散射并俘获
载流子，湿气腐蚀接触点，热斑使
二极管老化——这正是\omterm{def:b3:crystalline-solids:lattice}{晶体}在阳光下
服役三十年间缓慢积累的熵增。
\textbf{25.} 带隙收集了太阳光的一半；结区将
电子-空穴对转化为 \qty{36}{V} 的定向电流；串联组件、
遮挡与夏季高温各自分去一份；而
在 $\sim\qty{530}{kWh}$ 每年的发电量下，能量回收期为一年，
资金回收期为三年，裁决已定：装在屋顶上。
\end{solution}
